\documentclass[12pt,a4paper]{amsart}

\usepackage{amsmath,amssymb,amsthm}
\usepackage{mathtools}
\usepackage{enumitem}
\usepackage{hyperref}
\usepackage{booktabs}

%%─── Theorem environments ────────────────────────────────────────────
\newtheorem{theorem}{Theorem}[section]
\newtheorem{lemma}[theorem]{Lemma}
\newtheorem{proposition}[theorem]{Proposition}
\newtheorem{corollary}[theorem]{Corollary}
\newtheorem{conjecture}[theorem]{Conjecture}
\theoremstyle{definition}
\newtheorem{definition}[theorem]{Definition}
\newtheorem{problem}[theorem]{Open Problem}
\theoremstyle{remark}
\newtheorem{remark}[theorem]{Remark}
\newtheorem*{remark*}{Remark}

%%─── Macros ──────────────────────────────────────────────────────────
\newcommand{\ghat}{\widehat{g}}
\newcommand{\Phihat}{\widehat{\Phi}}
\newcommand{\Ai}{\mathrm{Ai}}
\newcommand{\Bi}{\mathrm{Bi}}
\newcommand{\ustar}{u_*}
\newcommand{\usaddle}{u^*}
\newcommand{\Fc}{F_c}
\newcommand{\norm}[1]{\|#1\|}
\newcommand{\abs}[1]{|#1|}
\newcommand{\ip}[2]{\langle #1,\, #2\rangle}
\DeclareMathOperator{\sgn}{sgn}
\newcommand{\OO}{\mathcal{O}}
\newcommand{\Zcross}{\mathcal{Z}_{\mathrm{cross}}}
\newcommand{\tstar}{t_*(c,n)}

%%─── AMS metadata ────────────────────────────────────────────────────
\subjclass[2020]{42C05, 34E20, 33E10, 47B06, 11M06}
\keywords{Prolate spheroidal wave functions, Mellin transform,
effective bandwidth, crossover zone, composite asymptotic,
Langer--Olver, stationary phase}

\begin{document}

\title[Crossover Zone and $c^{-1/2}$ Upper Bound for PSWFs]{%
Composite Asymptotics in the Crossover Zone
and a Rigorous $c^{-1/2}$ Upper Bound
for the Peak Width of Prolate Spheroidal
Wave Functions}

\author{[Author]}
\address{[Address]}
\email{[Email]}

\maketitle

\begin{abstract}
We continue the asymptotic analysis of the
Fourier--Mellin transforms $\widehat{\Phi}_n^{(c)}$
of the rescaled prolate spheroidal wave functions
(PSWFs) initiated in Papers~V--VI of this series.

The central obstacle identified in Paper~VI
\cite{PaperVI} was the \emph{crossover zone}
\[
\Zcross := \bigl\{t : c^{-2/3} \le |t - \tstar|
\le c^{-1/3}\bigr\},
\]
in which neither the Airy-profile approximation
nor the algebraic WKB bound applies uniformly.
We resolve this obstacle in two steps.

\textbf{Step~1 (Theorem~A).}
We construct a \emph{uniform composite
Airy--WKB asymptotic expansion} valid throughout
the Airy-to-WKB transition regime $\Zcross$:
\[
\widehat{\Phi}_n^{(c)}(t)
= c^{-1/2}\,\mathcal{A}_n^{(c)}\!\bigl(c^{1/2}(t-\tstar)\bigr)
\,\bigl(1 + O(c^{-1/3})\bigr),
\]
with
$|\mathcal{A}_n^{(c)}(x)| \le C_\kappa(1+|x|)^{-1/4}$
uniformly in $n \le \kappa c$.
The error $O(c^{-1/3})$ is controlled via
integration-by-parts remainder estimates
(the far-field remainder satisfies
$R_{\rm far} = O(c^{-5/6})$, giving a relative
error of order $c^{-1/3}$).

\textbf{Step~2 (Theorem~B).}
Using the composite asymptotic,
we derive the sharp upper bound
\[
\Delta_\varepsilon(c,n) \;\le\; C(\kappa,\varepsilon)\,c^{-1/2},
\]
improving the $c^{-2/3}$ upper bound of Paper~VI
and matching the $c^{-1/2}$ scaling observed
numerically in Paper~V \cite{PaperV}.

\textbf{The matching lower bound}
$\Delta_\varepsilon \ge c_1(\kappa,\varepsilon)\,c^{-1/2}$
is \emph{not} established here.
We show that all global $L^2$ mass-budget arguments
necessarily fail to produce the correct $c^{-1/2}$
lower bound, because the dominant $L^2$-mass of
$\widehat{\Phi}_n^{(c)}$ resides in the algebraic zone
$|t - \tstar| \ge c^{-1/3}$.
Any proof of the lower bound requires
\emph{pointwise non-cancellation control}
of $\widehat{\Phi}_n^{(c)}$ in a set of measure
$\gtrsim c^{-1/2}$ near $t_*$;
this is formulated precisely as
Problem~\ref{prob:nonvanish}
and is the central open problem for Paper~VIII.
\end{abstract}

%%═══════════════════════════════════════════════════════════════════
\section*{Summary of Main Results}
%%═══════════════════════════════════════════════════════════════════

\paragraph{Mass distribution (corrected).}
A key correction to Paper~VI, Remark~5.3:
the $L^2$-mass in the crossover zone is
$\asymp c^{-17/12}$, not $c^{-5/6}$ as stated there.
The corrected zone table is:

\begin{center}
\renewcommand{\arraystretch}{1.4}
\begin{tabular}{llll}
\toprule
Zone & $t$-range relative to $t_*$
& $|\widehat{\Phi}_n^{(c)}(t)|$
& $L^2$-mass \\
\midrule
Airy & $|t-t_*| \lesssim c^{-2/3}$
& $\asymp c^{-1/2}$
& $\sim c^{-5/3}$ \\
Crossover $\Zcross$
& $c^{-2/3} \le |t-t_*| \le c^{-1/3}$
& $\asymp c^{-1/2}|t-t_*|^{-1/4}$
& $\asymp c^{-17/12}$ \\
Algebraic & $|t-t_*| \ge c^{-1/3}$
& $\le C_\kappa c^{-1/4}t^{-1/4}$
& $\sim O(1)$ \textbf{(bulk)} \\
Barrier & $t \ge 2t_*$
& $\le e^{-\gamma_\kappa c^{1/2}}$
& negligible \\
\bottomrule
\end{tabular}
\end{center}

\noindent\textbf{Consequence.}
The bulk of the $L^2$-mass (the full $\tfrac{1}{2}$)
resides in the \emph{algebraic zone}.
Any lower bound on $\Delta_\varepsilon$
that relies only on global mass estimates
will overshoot to $\Delta \gtrsim c$
rather than $\Delta \gtrsim c^{-1/2}$
(Proposition~\ref{prop:lower-barrier}).

\paragraph{Chain of rigorous results.}
\[
\underbrace{\text{Composite in }\Zcross}
_{\text{Thm.~A}}
\;\Longrightarrow\;
\underbrace{|\widehat{\Phi}_n^{(c)}(t)|
\le C_\kappa c^{-1/2}(1+c^{1/2}|t-t_*|)^{-1/4}}
_{\text{Cor.~\ref{cor:uniform-decay}}}
\;\Longrightarrow\;
\underbrace{\Delta_\varepsilon \le C(\kappa,\varepsilon)\,c^{-1/2}}
_{\text{Thm.~B}}.
\]


%%─── Guide for the Referee ───────────────────────────────────────────
\section*{Guide for the Referee}

This paper has a single main thread:
establish the $c^{-1/2}$ \emph{upper bound} on $\Delta_\varepsilon$
by controlling the crossover zone.
The lower bound is not established and the reasons
why it is genuinely harder are explained precisely
in Section~\ref{sec:barrier}.

The logical structure is:
\begin{enumerate}
\item Section~\ref{sec:setup}: notation and
Paper~VI results used.
\item Section~\ref{sec:mass}: corrected mass
upper bound in $\Zcross$
(Proposition~\ref{prop:cross-mass};
proof uses only (VI.2), independent of Theorem~A),
followed by the mass partition
(Lemma~\ref{lem:mass-partition})
and the corrected zone table
(Corollary~\ref{cor:zone-table}).
\item Section~\ref{sec:composite}: construction
and rigorous justification of the composite
asymptotic (Theorem~A);
followed by the matching lower bound on the
crossover mass (Corollary~\ref{cor:cross-mass-lower}).
\item Section~\ref{sec:upper}: proof of the
$c^{-1/2}$ upper bound (Theorem~B).
\item Section~\ref{sec:barrier}: structural
analysis showing why the lower bound is
a different problem (Proposition~\ref{prop:lower-barrier}).
\item Section~\ref{sec:problems}: open problems.
\end{enumerate}

\tableofcontents

%%═══════════════════════════════════════════════════════════════════
\section{Setup and Reminders from Paper~VI}
\label{sec:setup}
%%═══════════════════════════════════════════════════════════════════

We retain all notation from Papers~V--VI.
Recall:
\begin{itemize}
\item $V_c(u) = c^2 e^{4u}
- (\chi_n^{(c)} + \tfrac{3}{4})\,e^{2u}
+ \tfrac{1}{4}$,
the potential in Schr\"odinger form
(Paper~VI, eq.~(2.2);
note the sign convention
$-\Phi'' + V_c \Phi = 0$,
so $V_c > 0$ in the classically
forbidden region $u \ll \ustar$).

\begin{remark}[Sign convention for $V_c$]
\label{rem:sign}
Throughout this paper we use the
Schr\"odinger form
$-(\Phi_n^{(c)})'' + V_c(u)\Phi_n^{(c)} = 0$
with
\[
V_c(u) = c^2 e^{4u}
- \bigl(\chi_n^{(c)} + \tfrac{3}{4}\bigr)e^{2u}
+ \tfrac{1}{4}.
\]
The turning point $\ustar$ is the unique root
of $V_c$ in $(-\infty, -\tfrac{1}{2}\log c)$,
satisfying $V_c(u) < 0$ for $u < \ustar$
(classically allowed region)
and $V_c(u) > 0$ for $u > \ustar$
(forbidden region).
The WKB momentum is
$p(u;c) = \sqrt{-V_c(u)}$
for $u < \ustar$
and $p(u;c) = \sqrt{V_c(u)}$
for $u > \ustar$.
This is consistent with
Papers~V and VI throughout.
\end{remark}
\item $\ustar = \ustar(c,n)$, unique turning point
of $V_c$ in $(-\infty, -\tfrac12\log c)$.
\item $p(u;c) = \sqrt{V_c(u)^+}$, WKB momentum.
\item $\tstar = p(\ustar;c) \asymp c^{1/2}$,
the saddle frequency.
\item $\usaddle = \usaddle(t;c,n)$, saddle point
defined by $p(\usaddle;c) = t$, for
$0 < t < \tstar$.
\item $\zeta(u;c,n)$, Langer variable,
$\tfrac{2}{3}|\zeta|^{3/2}
= |\int_{\ustar}^u p\,ds|$.
\item $D_n^{(c)}$, normalisation constant,
$|D_n^{(c)}| \asymp c^{1/4}$
(Paper~VI, Theorem~2.1).
\end{itemize}

The following results from Paper~VI are
used throughout without reproof:

\begin{enumerate}[(VI.1)]
\item\label{VI1} \textbf{Uniform Airy approximation}
(Thm.~2.1): For all $u \in (-\infty, 0)$,
\[
\Phi_n^{(c)}(u)
= \frac{D_n^{(c)}}{p(u;c)^{1/6}}
\Ai(\zeta(u;c,n))\,\bigl(1 + O(c^{-2/3})\bigr).
\]
\item\label{VI2} \textbf{Pointwise Mellin decay}
(Thm.~3.1): For $1 \le t \le c$,
\[
|\widehat{\Phi}_n^{(c)}(t)|
\le C_\kappa\,c^{-1/4}\,t^{-1/4}.
\]
\item\label{VI3} \textbf{Log-free $L^2$ tail}
(Thm.~3.2): For $1 \le R \le c^{1/4}$,
\[
\|\widehat{\Phi}_n^{(c)}\|_{L^2(|t|>R)}^2
\le C_\kappa^2\,c^{-1/2}\,R^{-1/2}.
\]
\item\label{VI4} \textbf{Fourier--Airy profile}
(Lem.~4.2): For $|t - \tstar| \le c^{-1/3}$,
\[
\widehat{\Phi}_n^{(c)}(t)
= c^{-1/2}\,\kappa_n^{(c)}\,
\Ai\!\bigl(c^{2/3}(t-\tstar)\bigr)
\,(1 + O(c^{-1/3})),
\quad
|\kappa_n^{(c)}| \asymp 1.
\]
\item\label{VI5} \textbf{$L^2$ normalisation}
(eq.~(2.4)):
$\|\widehat{\Phi}_n^{(c)}\|_{L^2(\mathbb{R})}^2
= \tfrac{1}{2}$.
\end{enumerate}

\begin{definition}[Crossover zone]
\label{def:Zcross}
\[
\Zcross(c,n)
:= \bigl\{t \in \mathbb{R} :
c^{-2/3} \le |t - \tstar| \le c^{-1/3}\bigr\}.
\]
In the rescaled variable $x := c^{1/2}(t - \tstar)$,
this becomes $\{x : c^{-1/6} \le |x| \le c^{1/6}\}$.
\end{definition}

\begin{remark}[Why the crossover zone is difficult]
\label{rem:difficulty}
The Airy profile \eqref{VI4} is valid for
$|c^{2/3}(t-\tstar)| = O(1)$,
i.e.\ $|t-\tstar| = O(c^{-2/3})$.
The WKB formula \eqref{VI2} requires the saddle
$\usaddle$ to lie well below $\ustar$, i.e.\
$|t - \tstar| \gg c^{-1/2}$ (so that
$\ustar - \usaddle \gg c^{-1/2}$).
In the crossover zone $\Zcross$, the gap
$\ustar - \usaddle \asymp c^{-1/6}$
is too large for the Airy approximation
and too small for the pure WKB formula.
Neither approximation applies, and a
\emph{composite} treatment is necessary.
\end{remark}

\begin{remark}[Three scales and the composite bridge]
\label{rem:scales}
Three natural rescalings appear in the analysis:
\begin{itemize}
\item \textbf{Airy scale:} $\xi = c^{2/3}(t - \tstar)$,
  natural for $|t - \tstar| \lesssim c^{-2/3}$
  (Airy profile \eqref{VI4}).
\item \textbf{WKB scale:} $y = c^{1/2}(t - \tstar)$,
  natural for $|t - \tstar| \gg c^{-1/2}$
  (algebraic decay \eqref{VI2}).
\item \textbf{Composite/Langer scale:}
  $x = c^{1/2}(t - \tstar)$ (the same scaling as the WKB variable),
  but with profile function $\mathcal{A}_n^{(c)}(x)$
  that reduces to the Airy function for
  $|x| \ll c^{1/6}$ via
  $\xi = c^{-1/6}\, x\,/\,(\mu_n^{(c)})^{2/3}$.
\end{itemize}
The Langer--Olver map (Lemma~\ref{lem:phase})
provides the explicit bridge:
$\mathcal{A}_n^{(c)}(x)
= \kappa_n^{(c)}\,\Ai\!\bigl(x/(\mu_n^{(c)})^{2/3}\bigr)$
interpolates between the Airy regime
($|x| \ll c^{1/6}$, Airy variable $\xi = c^{2/3}(t-\tstar)$)
and the WKB regime
($|x| \gg 1$, where $|\mathcal{A}| \lesssim |x|^{-1/4}$).
The composite variable $x$ is not ad hoc:
it is precisely the output of the Langer substitution
applied to the phase in Lemma~\ref{lem:phase}.
\end{remark}

%%═══════════════════════════════════════════════════════════════════
\section{Corrected Mass Upper Bound in $\Zcross$}
\label{sec:mass}
%%═══════════════════════════════════════════════════════════════════

\begin{proposition}[Upper bound on mass in $\Zcross$]
\label{prop:cross-mass}
For all $c \ge c_0(\kappa)$ and $n \le \kappa c$:
\begin{equation}
\label{eq:cross-mass-upper}
\int_{\Zcross}
|\widehat{\Phi}_n^{(c)}(t)|^2\,dt
\;\le\; C_\kappa\, c^{-17/12}.
\end{equation}
In particular, $c^{-17/12} \ll 1$,
so the crossover zone carries negligible $L^2$-mass
compared to the algebraic zone.
\end{proposition}

\begin{proof}
This proof uses only \eqref{VI2}; no appeal
to Theorem~A is made.

From \eqref{VI2}: $|\widehat{\Phi}_n^{(c)}(t)|
\le C_\kappa c^{-1/4} t^{-1/4}$.
For $t \in \Zcross$ we have
$t \asymp \tstar \asymp c^{1/2}$,
so $|\widehat{\Phi}_n^{(c)}(t)|^2
\le C_\kappa^2\, c^{-1/2}\, t^{-1/2}$.

Write $t = \tstar + r$ with
$c^{-2/3} \le |r| \le c^{-1/3}$.
Since $t \ge |r|$ we have $t^{-1/2} \le |r|^{-1/2}$.
Using the substitution $t = \tstar + c^{-1/2}x$,
$dt = c^{-1/2}\,dx$, and the refined amplitude
$|\widehat{\Phi}_n^{(c)}|^2 \le C_\kappa^2 c^{-1} |x|^{-1/2}$
(from \eqref{VI2} with $t \asymp c^{1/2}$):
\begin{align}
\int_{\Zcross} |\widehat{\Phi}_n^{(c)}(t)|^2\,dt
&\le C_\kappa^2\, c^{-1} \cdot c^{-1/2}
\int_{c^{-1/6}}^{c^{1/6}} x^{-1/2}\,dx
\notag\\
&= C_\kappa^2\, c^{-3/2}\cdot
\bigl[2x^{1/2}\bigr]_{c^{-1/6}}^{c^{1/6}}
= C_\kappa^2\,c^{-3/2}\cdot 2(c^{1/12}-c^{-1/12})
\notag\\
&\le 2C_\kappa^2\, c^{-3/2+1/12}
= 2C_\kappa^2\, c^{-17/12}.
\label{eq:cross-sharp}
\end{align}
\end{proof}

\begin{remark}[Matching lower bound deferred]
\label{rem:lower-deferred}
The matching lower bound
$\int_{\Zcross}|\widehat{\Phi}_n^{(c)}|^2\,dt
\ge c_\kappa\, c^{-17/12}$
requires a pointwise lower bound on
$|\mathcal{A}_n^{(c)}(x)|$,
which is available only after Theorem~A;
see Corollary~\ref{cor:cross-mass-lower}.
\end{remark}

\begin{lemma}[Mass partition]
\label{lem:mass-partition}
\begin{equation}
\label{eq:mass-partition}
\|\widehat{\Phi}_n^{(c)}\|_{L^2(\mathrm{alg})}^2
= \tfrac{1}{2}
- \|\widehat{\Phi}_n^{(c)}\|_{L^2(\mathrm{Airy})}^2
- \|\widehat{\Phi}_n^{(c)}\|_{L^2(\Zcross)}^2
= \tfrac{1}{2} - O(c^{-17/12}).
\end{equation}
\end{lemma}

\begin{proof}
This follows directly from \eqref{VI5}
(total $L^2$-mass $= \tfrac{1}{2}$),
Paper~VI, Cor.~4.3 (Airy mass $\sim c^{-5/3}$),
and the upper bound \eqref{eq:cross-mass-upper}
($\le C_\kappa c^{-17/12}$).
Since $c^{-5/3} \ll c^{-17/12}$
(as $\tfrac{5}{3} = \tfrac{20}{12} > \tfrac{17}{12}$),
the dominant correction is $O(c^{-17/12})$.
\end{proof}

\begin{corollary}[Corrected zone table]
\label{cor:zone-table}
The $L^2$-mass of $\widehat{\Phi}_n^{(c)}$
in each frequency zone satisfies:
\begin{align}
\label{eq:mass-Airy}
\|\widehat{\Phi}_n^{(c)}\|_{L^2(\mathrm{Airy})}^2
&\sim c^{-5/3}, \\
\label{eq:mass-cross}
\|\widehat{\Phi}_n^{(c)}\|_{L^2(\Zcross)}^2
&\asymp c^{-17/12}, \\
\label{eq:mass-alg}
\|\widehat{\Phi}_n^{(c)}\|_{L^2(\mathrm{alg})}^2
&= \tfrac{1}{2} - O(c^{-17/12}).
\end{align}
In particular, the algebraic zone carries
the full $L^2$-bulk.
\end{corollary}

\begin{proof}
Equation~\eqref{eq:mass-Airy}: Paper~VI, Cor.~4.3.

Upper bound in \eqref{eq:mass-cross}:
Proposition~\ref{prop:cross-mass}
(proved in this section using only \eqref{VI2},
independently of Theorem~A).

Lower bound in \eqref{eq:mass-cross}:
Corollary~\ref{cor:cross-mass-lower}
(Section~\ref{sec:composite}), which is proved
after Theorem~A and uses only
Theorem~A\eqref{thm:A-lower} and $|\kappa_n^{(c)}| \asymp 1$.
There is no circularity: the upper bound in
\eqref{eq:mass-cross} is used in the proof of
Theorem~A only via Lemma~\ref{lem:remainder},
which relies on \eqref{VI2} alone.
The lower bound in \eqref{eq:mass-cross}
is a consequence of Theorem~A, not an input to it.

Equation~\eqref{eq:mass-alg}:
Lemma~\ref{lem:mass-partition}.
\end{proof}

%%═══════════════════════════════════════════════════════════════════
\section{Composite Asymptotic Expansion}
\label{sec:composite}
%%═══════════════════════════════════════════════════════════════════

\subsection{Main statement (Theorem~A)}

\begin{theorem}[Composite Airy--WKB asymptotic]
\label{thm:A}
Fix $\kappa \in (0,2/\pi)$.
There exist constants $C_\kappa, c_\kappa, c_0(\kappa) > 0$
and, for each $c \ge c_0$ and $n \le \kappa c$,
a function $\mathcal{A}_n^{(c)} : \mathbb{R} \to \mathbb{C}$
such that for all $t \in \Zcross$:
\begin{equation}
\label{eq:composite-A}
\widehat{\Phi}_n^{(c)}(t)
= c^{-1/2}\,\mathcal{A}_n^{(c)}\!\bigl(c^{1/2}(t-\tstar)\bigr)
\;\bigl(1 + O(c^{-1/3})\bigr),
\end{equation}
where the $O$-constant is uniform in $n \le \kappa c$
and $t \in \Zcross$, and $\mathcal{A}_n^{(c)}$ satisfies:
\begin{enumerate}[\rm(a)]
\item\label{thm:A-upper}
\textbf{Upper bound:}
$|\mathcal{A}_n^{(c)}(x)|
\le C_\kappa\,(1+|x|)^{-1/4}$
for all $x \in \mathbb{R}$.
\item\label{thm:A-lower}
\textbf{Lower bound on intervals of length $\pi$:}
For every $x_0 \in \mathbb{R}$ there exists
$x_1 \in [x_0, x_0 + \pi]$ such that
$|\mathcal{A}_n^{(c)}(x_1)| \ge c_\kappa(1+|x_0|)^{-1/4}$.
\item\label{thm:A-match-airy}
\textbf{Matching to Airy zone:}
For $|x| \le c^{1/12}$,
$\mathcal{A}_n^{(c)}(x)
= \kappa_n^{(c)}\,\Ai\!\bigl(x/(\mu_n^{(c)})^{2/3}\bigr)
\,(1 + O(c^{-1/3}))$,
consistent with \eqref{VI4},
with $|\kappa_n^{(c)}| \asymp 1$.
\item\label{thm:A-match-WKB}
\textbf{Matching to algebraic zone:}
For $|x| \ge c^{1/12}$,
$|\mathcal{A}_n^{(c)}(x)| \le C_\kappa\,|x|^{-1/4}$,
consistent with \eqref{VI2}.
\end{enumerate}
\end{theorem}

\begin{remark}[Explicit form and well-definedness
of $\kappa_n^{(c)}$]
\label{rem:A-explicit}
The profile function is
\begin{equation}
\label{eq:A-explicit}
\mathcal{A}_n^{(c)}(x)
:= \kappa_n^{(c)}\,
\Ai\!\left(\frac{x}{(\mu_n^{(c)})^{2/3}}\right),
\qquad |\kappa_n^{(c)}| \asymp 1,
\end{equation}
where
\begin{equation}
\label{eq:mu-def}
\mu_n^{(c)} := p''(\ustar;c)\,c^{-1/2}
\;\in\; [c_\kappa, C_\kappa]
\end{equation}
is the rescaled cubic-phase coefficient
(Lemma~\ref{lem:uniform-n}).

\textbf{Definition of $\kappa_n^{(c)}$.}
The coefficient $\kappa_n^{(c)}$ is defined
by matching \eqref{eq:composite-A} with
the Fourier--Airy profile \eqref{VI4} at $t = \tstar$
(i.e.\ $x = 0$).
At $x = 0$, equation \eqref{VI4} gives:
\[
\widehat{\Phi}_n^{(c)}(\tstar)
= c^{-1/2}\,\kappa_n^{(c,\mathrm{VI})}\,\Ai(0)
\,(1 + O(c^{-1/3})),
\]
where $\kappa_n^{(c,\mathrm{VI})}$ is the
constant from \eqref{VI4} with
$|\kappa_n^{(c,\mathrm{VI})}| \asymp 1$.
From \eqref{eq:A-explicit} at $x = 0$:
\[
c^{-1/2}\mathcal{A}_n^{(c)}(0)
= c^{-1/2}\kappa_n^{(c)}\,\Ai(0).
\]
Matching these two expressions defines:
\begin{equation}
\label{eq:kappa-def}
\kappa_n^{(c)}
:= \kappa_n^{(c,\mathrm{VI})}
\,(1 + O(c^{-1/3})).
\end{equation}

\textbf{Well-definedness and precision.}
Since \eqref{VI4} holds with error $O(c^{-1/3})$,
the coefficient $\kappa_n^{(c)}$ is determined
only modulo $O(c^{-1/3})$:
\[
\kappa_n^{(c)}
= \kappa_n^{(c,\mathrm{VI})} + O(c^{-1/3}).
\]
This is consistent with the error term
$O(c^{-1/3})$ in \eqref{eq:composite-A}:
replacing $\kappa_n^{(c)}$ by
$\kappa_n^{(c)} + O(c^{-1/3})$
changes \eqref{eq:composite-A} by
$O(c^{-1/2} \cdot c^{-1/3}) = O(c^{-5/6})$,
which is within the stated error.
No sharper determination of $\kappa_n^{(c)}$
is needed for the results of this paper.

\textbf{Lower bound $|\kappa_n^{(c)}| \asymp 1$.}
From \eqref{eq:kappa-def} and
$|\kappa_n^{(c,\mathrm{VI})}| \asymp 1$
(Paper~VI, Lem.~4.2), we obtain
$|\kappa_n^{(c)}| \ge c_\kappa > 0$
for all $c \ge c_0(\kappa)$ and $n \le \kappa c$.
This lower bound is used in
Theorem~A\eqref{thm:A-lower}
and Corollary~\ref{cor:cross-mass-lower}.

\textbf{Dependence on $\mu_n^{(c)}$.}
The rescaling of the Airy argument by
$(\mu_n^{(c)})^{2/3}$ arises from the
cubic-phase normalisation in the proof of
Theorem~A (substitution $s' = (\mu_n^{(c)})^{-1/3}y$).
Since $\mu_n^{(c)} \in [c_\kappa, C_\kappa]$
uniformly (Lemma~\ref{lem:uniform-n}),
the Airy function $\Ai(x/(\mu_n^{(c)})^{2/3})$
is uniformly equivalent to $\Ai(x/C)$
for a fixed constant $C \in [c_\kappa^{2/3},
C_\kappa^{2/3}]$,
and all estimates in Theorem~A\eqref{thm:A-upper}
and \eqref{thm:A-lower} hold with
constants depending only on $\kappa$.
\end{remark}

\begin{remark}[Composite nature of $\mathcal{A}_n^{(c)}$]
\label{rem:composite-nature}
Although the leading term in \eqref{eq:A-explicit}
is an Airy function, Theorem~A is \emph{not} a
restatement of the Airy profile \eqref{VI4}.
The composite nature lies in three features
absent from \eqref{VI4}:
\begin{enumerate}[\rm(i)]
\item \textbf{Uniform validity:}
\eqref{eq:composite-A} holds throughout $\Zcross$,
where $|c^{2/3}(t-\tstar)|$ ranges up to $c^{1/3} \to \infty$
and the Airy profile \eqref{VI4} breaks down.
\item \textbf{Remainder control:}
the error $O(c^{-1/3})$ is rigorously bounded
via the integration-by-parts estimate of
Lemma~\ref{lem:remainder}, not assumed.
\item \textbf{Domain extension:}
the local integral over $|s'| \le c^{1/6}$
is extended to $\mathbb{R}$ with controlled error;
this extension is what makes \eqref{eq:composite-A}
a genuine asymptotic expansion rather than
a local approximation.
\end{enumerate}
\end{remark}

\subsection{Proof of Theorem~A}

The proof consists of four lemmas.

\begin{lemma}[Phase expansion and Airy normal form]
\label{lem:phase}
Write $t = \tstar + c^{-1/2}x$,
$u = \ustar + c^{-1/2}s$.
For $|s|, |x| \le c^{1/6}$:
\[
c^{1/2}[\phi_+(u;t,c) - \phi_0]
= xs + \tfrac{\mu_n^{(c)}}{6}s^3 + O(c^{-1/3}),
\]
uniformly, where $\phi_0 = S(\ustar;c)+\tstar\ustar$.

Moreover, the phase is in Airy normal form
up to admissible error: the quadratic term,
though non-negligible for $|s| \asymp c^{1/6}$
(of order $c^{1/12}$ after rescaling),
is absorbed by a smooth change of variables
$s \mapsto s'(s,x)$ depending on $x$,
as in Olver~\cite{Olver1974}, Ch.~11.
This transformation yields a phase of the form
$x's' + \tfrac{\mu_n^{(c)}}{6}s'^3 + O(c^{-1/3})$
uniformly for $|x| \le c^{1/6}$,
with $s' = s + O(c^{-1/6})$,
so the integration domain is preserved
up to admissible errors.
\end{lemma}

\begin{proof}
Taylor-expand $S(u;c)$ around $\ustar$;
linear terms cancel by the saddle condition
$p(\ustar;c) = \tstar$.
Quadratic: $\tfrac{1}{2}p'(\ustar)c^{-1}s^2
\asymp c^{-3/4}s^2$; after $c^{1/2}$-rescaling
this is $c^{-1/4}s^2$, which for $s \asymp c^{1/6}$
gives $c^{1/12}$, hence \emph{not} negligible.
The Olver normal-form reduction
(\cite{Olver1974}, Ch.~11, Theorem~11.1)
elaborates a smooth diffeomorphism $s \mapsto s'(s,x)$
that eliminates the quadratic term,
yielding the pure cubic phase structure.
Cubic term: $\tfrac{1}{6}\mu_n^{(c)}s'^3$
by \eqref{eq:mu-def}.
Quartic and higher: $O(c^{-1/3})$
after the normal-form substitution.

The diffeomorphism satisfies the following
uniform estimates for $|s|, |x| \le c^{1/6}$:
\begin{enumerate}[\rm(i)]
\item \textbf{Uniformity in $n$:}
$s'(s,x)$ depends only on $\mu_n^{(c)} \in [c_\kappa, C_\kappa]$
(Lemma~\ref{lem:uniform-n}),
so all constants in the normal-form reduction
are uniform in $n \le \kappa c$.
\item \textbf{Jacobian control:}
$\dfrac{\partial s'}{\partial s} = 1 + O(c^{-1/6})$,
so $ds' = (1 + O(c^{-1/6}))\,ds$;
the Jacobian contributes only to the
$O(c^{-1/3})$ error in \eqref{eq:composite-A}.
\item \textbf{Domain preservation:}
For $|s| \le c^{1/6}$ one has $|s'| \le 2c^{1/6}$,
so the integration domain $|s'| \le c^{1/6}$
(enlarged by the factor $2$ to absorb the shift
$s' = s + O(c^{-1/6})$)
incurs no additional boundary loss beyond
the tail already controlled by
Lemma~\ref{lem:remainder}.
\item \textbf{Explicit diffeomorphism and uniform error
on the growing window.}
We construct $s \mapsto s'(s,x)$ explicitly to second
order and verify that the Olver remainder is
$O(c^{-1/3})$ uniformly on $|s|, |x| \le c^{1/6}$.

\textbf{Step~1: Taylor expansion of the phase.}
Write $u = \ustar + c^{-1/2}s$,
$t = \tstar + c^{-1/2}x$.
The rescaled oscillatory phase becomes:
\[
\Psi(s,x) := c^{1/2}\phi_+(u;t,c)
= -xs + \tfrac{\alpha_2}{2}s^2
+ \tfrac{\mu_n^{(c)}}{6}s^3
+ \tfrac{\alpha_4}{24}s^4 c^{-1/2} + \cdots,
\]
where
\[
\alpha_2 := p'(\ustar)\,c^{-1/2} = O(c^{-1/4}),
\quad
\mu_n^{(c)} := p''(\ustar)\,c^{-1/2}
\in [c_\kappa, C_\kappa],
\quad
\alpha_4 := p'''(\ustar)\,c^{-1/2} = O(1).
\]

\textbf{Step~2: Elimination of the quadratic term.}
The quadratic elimination requires
\[
s' = s + \frac{\alpha_2}{\mu_n^{(c)}} + O(s^2\alpha_2),
\]
so the shift is $\beta = \alpha_2/\mu_n^{(c)} = O(c^{-1/4})$.
After substitution the phase becomes:
\[
\Psi = -x's' + \tfrac{\mu_n^{(c)}}{6}s'^3
+ R_4(s', x),
\]
where $x' = x + O(\alpha_2)$ and the quartic remainder satisfies
$|R_4(s',x)| \le C_\kappa\,c^{-1/2}\,|s'|^4$
for $|s'| \le c^{1/6}$.

\textbf{Step~3: Uniform error estimate.}
After undoing the $c^{1/2}$-rescaling:
\[
|c^{-1/2} R_4(s',x)|
\le C_\kappa\,c^{-1}\,|s'|^4
\le C_\kappa\,c^{-1}\,(c^{1/6})^4
= C_\kappa\,c^{-1/3}
\quad \text{for } |s'| \le c^{1/6}.
\]
This is uniform over the entire growing window.

\textbf{Step~4: Multiplicative error.}
The phase error $O(c^{-1/3})$ gives a multiplicative error:
\[
e^{i(\Psi + c^{-1/3}\theta)}
= e^{i\Psi}(1 + O(c^{-1/3}))
\]
uniformly for $|\theta| \le C_\kappa$,
so the Olver normal-form reduction is valid
with relative error $O(c^{-1/3})$
uniformly on the growing window $|s'| \le c^{1/6}$
for all $c \ge c_0(\kappa)$ and $n \le \kappa c$.

\textbf{Step~5: Jacobian.}
The Jacobian satisfies
$\partial s'/\partial s = 1 + O(c^{-1/4})$
for $|s| \le c^{1/6}$,
confirming part~(ii) above.
\end{enumerate}
\end{proof}

\begin{lemma}[Amplitude near $\ustar$]
\label{lem:amp-near}
For $|u - \ustar| \le c^{-1/3}$:
$\Phi_n^{(c)}(u)
= D_n^{(c)}\tstar^{-1/6}\Ai(0)
(1 + O(c^{-1/12}))$.
\end{lemma}

\begin{proof}
From \eqref{VI1}: $|\zeta(u)| \le C_\kappa c^{-1/6}\to 0$
so $\Ai(\zeta)=\Ai(0)+O(c^{-1/6})$;
$|p(u)-\tstar| \le C_\kappa c^{-1/12}$
(since $p'(\ustar)\asymp c^{1/4}$)
so $p(u)^{-1/6} = \tstar^{-1/6}(1+O(c^{-1/12}))$.

For the derivative bound used in Lemma~\ref{lem:remainder},
differentiate \eqref{VI1} with respect to $u$:
\[
(\Phi_n^{(c)})'(u)
= D_n^{(c)}\Bigl[
-\tfrac{1}{6}\frac{p'(u)}{p(u)^{7/6}}\Ai(\zeta)
+ \frac{\zeta'(u)}{p(u)^{1/6}}\Ai'(\zeta)
\Bigr]\bigl(1 + O(c^{-2/3})\bigr).
\]
For $|u - \ustar| \le c^{-1/3}$:
$|p(u)| \asymp \tstar \asymp c^{1/2}$,
$|p'(u)| \asymp c^{1/4}$,
$|\zeta'(u)| = p(u) \asymp c^{1/2}$,
and $|\Ai(\zeta)|, |\Ai'(\zeta)| = O(1)$.
Using $|D_n^{(c)}| \asymp c^{1/4}$:
\[
|(\Phi_n^{(c)})'(u)|
\le C_\kappa c^{1/4}
\bigl(c^{1/4} \cdot c^{-7/12} + c^{1/2} \cdot c^{-1/12}\bigr)
= C_\kappa c^{1/4} \cdot c^{5/12}
\asymp C_\kappa\, c^{2/3}.
\]
After the rescaling $s = c^{1/2}(u - \ustar)$,
so $a(s) = c^{-1/2}\Phi_n^{(c)}(\ustar + c^{-1/2}s)$
and $a'(s) = c^{-1}(\Phi_n^{(c)})'(\ustar + c^{-1/2}s)$:
\[
|a'(s)| \le c^{-1} \cdot C_\kappa c^{2/3}
= C_\kappa\, c^{-1/3}.
\]


\end{proof}

\begin{lemma}[Tail remainder]
\label{lem:remainder}
For $t \in \Zcross$:
\[
\Bigl|\int_{|u-\ustar|>c^{-1/3}}
\Phi_n^{(c)}(u)e^{itu}\,du\Bigr|
= O(c^{-5/6}).
\]
The relative error in \eqref{eq:composite-A}
is therefore $O(c^{-1/3})$.
\end{lemma}

\begin{proof}
All estimates below are performed entirely in the
rescaled variable $s = c^{1/2}(u-\ustar)$;
the Jacobian $du = c^{-1/2}\,ds$ is absorbed into
the amplitude $a(s)$, which satisfies
$\|a\|_{L^\infty} \le C_\kappa c^{-1/2}$
(combining the $c^{-1/2}$ from the substitution with
the $O(1)$ amplitude of $\Phi_n^{(c)}$ near $\ustar$,
uniformly in the rescaled regime near the turning point,
cf.\ Lemma~\ref{lem:amp-near}).
The main term is therefore $O(c^{-1/2})$,
and all error bounds below are stated as
fractions of this main term.

The phase is $\varphi(s) = xs + \tfrac{\mu_n^{(c)}}{6}s^3$
with derivative $\varphi'(s) = x + \tfrac{\mu_n^{(c)}}{2}s^2$.
For $|s| \ge c^{1/6}$ and $t \in \Zcross$
(so $|x| \le c^{1/6}$), the dominant term satisfies
$|\tfrac{\mu}{2}s^2| \asymp s^2 \ge c^{1/3}$,
hence $|\varphi'(s)| \asymp s^2 \ge c^{1/3}$.
Moreover, $|\varphi''(s)| = |\mu_n^{(c)} s| \asymp |s|$,
so
\[
\frac{|\varphi''(s)|}{|\varphi'(s)|^2}
\lesssim \frac{|s|}{s^4} = s^{-3}.
\]

Integration by parts yields
\[
\int_{|s|\ge c^{1/6}} e^{i\varphi(s)} a(s)\,ds
=
\left[\frac{e^{i\varphi(s)}a(s)}{i\varphi'(s)}\right]_{|s|=c^{1/6}}
-
\int_{|s|\ge c^{1/6}} e^{i\varphi(s)}
\frac{d}{ds}\!\left(\frac{a(s)}{\varphi'(s)}\right)\,ds.
\]

\noindent\textit{Boundary term.}
At $|s| = c^{1/6}$ we have $|\varphi'(s)| \gtrsim c^{1/3}$
and $|a(s)| = O(c^{-1/2})$, so
\[
\left|\frac{a(s)}{\varphi'(s)}\right|_{|s|=c^{1/6}}
\lesssim c^{-1/2} \cdot c^{-1/3} = c^{-5/6}.
\]

\noindent\textit{Remainder integral.}
Since $|a'(s)| = O(c^{-1/3})$ (Lemma~\ref{lem:amp-near})
and $|\varphi'(s)| \gtrsim s^2$,
\[
\left|\frac{d}{ds}\left(\frac{a(s)}{\varphi'(s)}\right)\right|
\lesssim \frac{|a'(s)|}{|\varphi'(s)|}
+ |a(s)| \frac{|\varphi''(s)|}{|\varphi'(s)|^2}
\lesssim c^{-1/3} s^{-2} + c^{-1/2} s^{-3}
\lesssim c^{-1/3} s^{-2}.
\]
Hence
\[
\int_{|s|\ge c^{1/6}}
\left|\frac{d}{ds}\!\left(\frac{a(s)}{\varphi'(s)}\right)\right|ds
\lesssim c^{-1/3} \int_{c^{1/6}}^{\infty} s^{-2}\,ds
\lesssim c^{-1/3} \cdot c^{-1/6} = c^{-1/2}.
\]

The boundary term is $O(c^{-5/6})$ absolutely.
The remainder integral is $O(c^{-1/2})$ absolutely,
matching the main-term order; hence it contributes
$O(c^{-1/3})$ \emph{relative} to the main term $c^{-1/2}$.
The total tail is therefore $O(c^{-1/3})$ relative,
as claimed in~\eqref{eq:composite-A}.
\end{proof}

\begin{lemma}[Uniformity in $n$]
\label{lem:uniform-n}
$c_\kappa \le \mu_n^{(c)} \le C_\kappa$
uniformly in $n \le \kappa c$, $c \ge c_0$.
\end{lemma}

\begin{proof}
At the turning point $\ustar$ we have $V_c(\ustar) = 0$,
so $p(u) = \sqrt{V_c(u)}$ satisfies
\[
p'(u) = \frac{V_c'(u)}{2p(u)}.
\]
Although $p'(u)$ is singular at $\ustar$,
the singularity is removable.
Applying L'H\^opital's rule as $u \to \ustar$ gives
\[
p'(\ustar)^2 = \tfrac{1}{2}V_c''(\ustar),
\]
and differentiating once more yields
\[
p''(\ustar) \asymp \bigl(V_c''(\ustar)\bigr)^{1/2}.
\]
Since $V_c(u) = c^2 e^{4u} - \chi_n^{(c)} e^{2u} + \tfrac{1}{4}$,
a direct computation gives
$V_c''(\ustar) = 16c^2 e^{4\ustar} - 4\chi_n^{(c)} e^{2\ustar}
= 4e^{2\ustar}(4c^2 e^{2\ustar} - \chi_n^{(c)})$.
Using $V_c(\ustar)=0$ to substitute
$\chi_n^{(c)} e^{2\ustar} = c^2 e^{4\ustar} + \tfrac{1}{4}$
one obtains $V_c''(\ustar) \asymp c^2 e^{4\ustar} \asymp c$
(since $e^{2\ustar} \asymp c^{-1/2}$ at the turning point).
Hence $p''(\ustar) \asymp c^{1/2}$,
so $\mu_n^{(c)} = p''(\ustar)\,c^{-1/2} \in [c_\kappa, C_\kappa]$
uniformly in $n \le \kappa c$.
\end{proof}

\begin{proof}[Proof of Theorem~A]
Substitute $u = \ustar + c^{-1/2}s$
and apply the Olver normal-form reduction
of Lemma~\ref{lem:phase}.
Use Lemmata~\ref{lem:amp-near}--\ref{lem:uniform-n}:
\[
\widehat{\Phi}_n^{(c)}(t)
= c^{-1/2}e^{i\phi_0}\kappa_n^{(c)}
\int_{|s'|\le c^{1/6}}
e^{i(x's' + \mu_n^{(c)}s'^3/6)}\,ds'
\cdot(1+O(c^{-1/12}))
+ R_{\rm far}.
\]
To evaluate the local integral, set
$s' = (\mu_n^{(c)})^{-1/3}y$;
then the phase becomes
$(\mu_n^{(c)})^{-1/3}x'y + y^3/3$,
and by the Airy integral identity
$\int_{-\infty}^\infty e^{i(\xi y + y^3/3)}\,dy
= 2\pi\,\Ai(\xi)$
with $\xi = (\mu_n^{(c)})^{-1/3}x'$:
\[
\int_{-\infty}^\infty
e^{i(x's' + \mu_n^{(c)}s'^3/6)}\,ds'
= 2\pi\,(\mu_n^{(c)})^{-1/3}\,
\Ai\!\Bigl(\frac{x'}{(\mu_n^{(c)})^{2/3}}\Bigr).
\]
Extending from $|s'| \le c^{1/6}$ to $\mathbb{R}$
incurs an error controlled by
Lemma~\ref{lem:remainder}.

\begin{remark*}[Remainder after normal-form transformation]
We verify that Lemma~\ref{lem:remainder} applies to the
transformed integral. After the Olver diffeomorphism
$s \mapsto s'(s,x)$ of Lemma~\ref{lem:phase}, the integrand
becomes $\tilde{a}(s')\,e^{i\tilde{\varphi}(s')}$, where
\[
\tilde{a}(s') := a(s(s',x))\,\frac{\partial s}{\partial s'},
\qquad
\tilde{\varphi}(s') := x's' + \tfrac{\mu_n^{(c)}}{6}s'^3.
\]
We check the two amplitude conditions used in
Lemma~\ref{lem:remainder}:

\textbf{(i) $L^\infty$-bound on $\tilde{a}$.}
Since $\|\tilde{a}\|_{L^\infty}
= \|a \cdot \partial_s s'\|_{L^\infty}
\le \|a\|_{L^\infty} \cdot \|\partial_s s'\|_{L^\infty}$,
and Lemma~\ref{lem:phase}(ii) gives
$\partial_s s' = 1 + O(c^{-1/6})$,
we obtain $\|\tilde{a}\|_{L^\infty}
= \|a\|_{L^\infty}(1 + O(c^{-1/6}))
= O(c^{-1/2})$, consistent with the bound
used in Lemma~\ref{lem:remainder}.

\textbf{(ii) Derivative bound on $\tilde{a}$.}
By the chain and product rules:
\[
\tilde{a}'(s') = a'(s(s',x))\,\Bigl(\frac{\partial s}{\partial s'}\Bigr)^2
+ a(s(s',x))\,\frac{\partial^2 s}{\partial s'^2}.
\]
Lemma~\ref{lem:phase}(ii) gives
$|\partial^2 s/\partial s'^2| = O(c^{-1/6})$,
so
\[
|\tilde{a}'(s')| \le |a'| \cdot (1+O(c^{-1/6}))^2
+ |a| \cdot O(c^{-1/6})
= O(c^{-1/3}) + O(c^{-1/6} \cdot c^{-1/2})
= O(c^{-1/3}),
\]
using $|a'| = O(c^{-1/3})$ and $|a| = O(c^{-1/2})$
from Lemma~\ref{lem:amp-near}.

\textbf{(iii) Phase derivative bound on $\tilde{\varphi}$.}
The transformed phase satisfies
$\tilde{\varphi}'(s') = x' + \tfrac{\mu_n^{(c)}}{2}s'^2$,
which has the same structure as
$\varphi'(s)$ in Lemma~\ref{lem:remainder},
with $|\tilde{\varphi}'(s')| \asymp s'^2$
for $|s'| \ge c^{1/6}$.

Since $\tilde{a}$ satisfies the same $L^\infty$-
and derivative bounds as $a$, and $\tilde{\varphi}$
the same phase-derivative lower bound as $\varphi$,
the argument of Lemma~\ref{lem:remainder} applies
verbatim to $\tilde{a}$ and $\tilde{\varphi}$,
yielding the same $O(c^{-5/6})$ boundary term
and $O(c^{-1/3})$ relative remainder.
\end{remark*}

Combining, and writing $x = c^{1/2}(t - \tstar)$:
\[
\widehat{\Phi}_n^{(c)}(t)
= c^{-1/2}\kappa_n^{(c)}
\Ai\!\Bigl(\frac{x}{(\mu_n^{(c)})^{2/3}}\Bigr)
(1+O(c^{-1/3})),
\]
which is \eqref{eq:composite-A}.
Properties~(a)--(d) follow from standard Airy bounds,
the oscillation of $\Ai$, and matching with \eqref{VI4}.
\end{proof}

\begin{corollary}[Lower bound on mass in $\Zcross$]
\label{cor:cross-mass-lower}
For all $c \ge c_0(\kappa)$, $n \le \kappa c$:
\[
\int_{\Zcross}|\widehat{\Phi}_n^{(c)}(t)|^2\,dt
\;\ge\; c_\kappa\,c^{-17/12}.
\]
Together with Proposition~\ref{prop:cross-mass}
this gives $\asymp c^{-17/12}$.
\end{corollary}

\begin{proof}
By Theorem~A\eqref{thm:A-lower},
for every $x_0 \in [c^{-1/6},c^{1/6}]$
there is $x_1 \in [x_0,x_0+\pi]$ with
$|\mathcal{A}_n^{(c)}(x_1)| \ge c_\kappa|x_0|^{-1/4}$.
Hence via \eqref{eq:composite-A}:
\[
\int_{\Zcross}|\widehat{\Phi}|^2\,dt
\ge c_\kappa^2 c^{-3/2}
\int_{c^{-1/6}}^{c^{1/6}} x^{-1/2}\,dx
= c_\kappa^2 c^{-3/2}\cdot 2(c^{1/12}-c^{-1/12})
\ge c_\kappa^2 c^{-17/12}.
\]
\end{proof}

%%═══════════════════════════════════════════════════════════════════
\section{The $c^{-1/2}$ Upper Bound}
\label{sec:upper}
%%═══════════════════════════════════════════════════════════════════

\begin{corollary}[Uniform decay in $\Zcross$]
\label{cor:uniform-decay}
For $t \in \Zcross$, $c \ge c_0$, $n \le \kappa c$:
\begin{equation}
\label{eq:uniform-decay}
|\widehat{\Phi}_n^{(c)}(t)|
\le C_\kappa c^{-1/2}(1+c^{1/2}|t-\tstar|)^{-1/4}.
\end{equation}
\end{corollary}

\begin{proof}
Immediate from Theorem~A\eqref{thm:A-upper}
and \eqref{eq:composite-A},
since $x = c^{1/2}(t-\tstar)$ gives
$(1+|x|)^{-1/4} = (1+c^{1/2}|t-\tstar|)^{-1/4}$.
\end{proof}

\begin{theorem}[Upper bound on peak width, Theorem~B]
\label{thm:B}
For every $\kappa\in(0,2/\pi)$ and $\varepsilon\in(0,1)$
there exists $C(\kappa,\varepsilon)>0$ such that
for all $c \ge c_0(\kappa,\varepsilon)$, $n \le \kappa c$:
\begin{equation}
\label{eq:thm-B}
\Delta_\varepsilon(c,n) \le C(\kappa,\varepsilon)\,c^{-1/2}.
\end{equation}
\end{theorem}

\begin{proof}
We estimate the $L^2$-mass of $\widehat{\Phi}_n^{(c)}$
outside $[t_* - \Delta, t_* + \Delta]$
for $\Delta = K c^{-1/2}$, $K = K(\kappa,\varepsilon)$
chosen below.

\textbf{Step~1: Decomposition.}
Write
$\|\widehat{\Phi}_n^{(c)}\|_{L^2([t_*-\Delta, t_*+\Delta]^c)}^2
= J_1 + J_2 + J_3$,
where $J_1$ integrates over $\Zcross \cap \{|t-t_*|>\Delta\}$,
$J_2$ over the algebraic zone
$c^{-1/3} < |t-t_*| \le t_*$,
and $J_3$ over $|t-t_*| > t_*$.
Since $\Delta = Kc^{-1/2} \gg c^{-2/3}$ for large $c$,
the Airy zone lies inside $[t_*-\Delta, t_*+\Delta]$
and contributes nothing to $J_1$.

\textbf{Step~2: Bound on $J_1$.}
$J_1 \le \|\widehat{\Phi}_n^{(c)}\|^2_{L^2(\Zcross)}
\le C_\kappa c^{-17/12}$
by Proposition~\ref{prop:cross-mass}.
For $c \ge c_1(\kappa,\varepsilon)$ we have
$C_\kappa c^{-17/12} \le \varepsilon^2/6$.

\textbf{Step~3: Bound on $J_3$.}
For $t > 2t_*$, the barrier estimate
(Paper~V, Thm.~3.1) gives exponential decay;
for $t < 0$, $\widehat{\Phi}_n^{(c)} = 0$.
Hence $J_3 \le e^{-\gamma_\kappa c^{1/2}} \le \varepsilon^2/6$
for $c \ge c_2(\kappa,\varepsilon)$.

\textbf{Step~4: Bound on $J_2$.}
For $t$ in the algebraic zone,
$t \asymp t_* \asymp c^{1/2}$, so
\eqref{VI2} gives
$|\widehat{\Phi}_n^{(c)}(t)|^2
\le C_\kappa^2 c^{-1/2} t^{-1/2}
\le C_\kappa^2 c^{-1/2} (t_*)^{-1/2}
\asymp C_\kappa^2 c^{-3/4}$.
The integration region in $J_2$ for
$|t - t_*| > \Delta$ has
$t \in [t_*+\Delta, 2t_*]$ or $t \in [0, t_*-\Delta]$.
In both sub-regions, using $t^{-1/2} \le \Delta^{-1/2}$
(since $t \ge \Delta$ and $t \ge t_* - t_* = 0$
trivially, but more usefully
$t_* - \Delta \ge \Delta$ for $K \le t_*c^{1/2}/2$,
which holds for large $c$):
\[
J_2 \le C_\kappa^2 c^{-1/2} \cdot \Delta^{-1/2}
\cdot 2t_*
= C_\kappa^2 c^{-1/2} \cdot (Kc^{-1/2})^{-1/2}
\cdot 2c^{1/2}
= 2C_\kappa^2 K^{-1/2}.
\]
Choose $K = (6C_\kappa^2/\varepsilon^2)^2$
so that $J_2 \le \varepsilon^2/3$.

\textbf{Conclusion.}
$J_1 + J_2 + J_3
\le \varepsilon^2/6 + \varepsilon^2/3 + \varepsilon^2/6
= \frac{2\varepsilon^2}{3} < \varepsilon^2$
for $c \ge \max(c_1, c_2)(\kappa,\varepsilon)$.
This shows $\Delta_\varepsilon(c,n) \le Kc^{-1/2}
= C(\kappa,\varepsilon) c^{-1/2}$.
\end{proof}

\begin{remark}[Explicit constant]
\label{rem:explicit-const}
The proof yields the explicit value
$C(\kappa,\varepsilon) = (8C_\kappa^2/\varepsilon^2)^2$,
where $C_\kappa$ is the constant from \eqref{VI2}.
This constant is not claimed to be optimal;
the sharp value is the subject of
Problem~\ref{prob:sharp-const}.
\end{remark}

%%═══════════════════════════════════════════════════════════════════
\section{Why the Lower Bound is a Different Problem}
\label{sec:barrier}
%%═══════════════════════════════════════════════════════════════════

\begin{proposition}[Global mass arguments
cannot give the $c^{-1/2}$ lower bound]
\label{prop:lower-barrier}
Let $\mathcal{M}$ be any argument that uses
only the following two inputs:
\begin{enumerate}[\rm(a)]
\item the total $L^2$-normalisation
$\|\widehat{\Phi}_n^{(c)}\|_{L^2}^2
= \tfrac{1}{2}$,
\item pointwise upper bounds of the form
$|\widehat{\Phi}_n^{(c)}(t)|
\le f(t,c)$ for some explicit $f$.
\end{enumerate}
Then $\mathcal{M}$ cannot establish
$\Delta_\varepsilon(c,n) \ge c_1\,c^{-1/2}$
for any $c_1 > 0$.

More precisely: the best lower bound
on $\Delta_\varepsilon$ that $\mathcal{M}$
can produce is
\[
\Delta_\varepsilon(c,n)
\;\ge\; c_1(\kappa,\varepsilon)\,c^{-1},
\]
which is weaker than $c^{-1/2}$
by a factor of $c^{1/2}$.
\end{proposition}

\begin{proof}
A lower bound on $\Delta_\varepsilon$
via a mass argument takes the following form:
suppose the $L^2$-mass on
$[t_* - \Delta, t_* + \Delta]$
satisfies
\[
\|\widehat{\Phi}_n^{(c)}\|^2_{L^2(|t-t_*|
\le \Delta)}
\le 2\Delta \cdot \sup_{|t-t_*|\le\Delta}
|\widehat{\Phi}_n^{(c)}(t)|^2
\le 2\Delta \cdot f(\Delta, c)^2.
\]
For the tail to exceed $\varepsilon^2\cdot\tfrac{1}{2}$
(i.e.\ the peak to contain less than
$(1-\varepsilon^2)$ fraction of the total mass),
one needs:
\[
2\Delta \cdot f(\Delta,c)^2
< (1-\varepsilon^2) \cdot \tfrac{1}{2},
\]
which gives a lower bound
$\Delta \ge c_1/f(\Delta,c)^2$.

\textbf{Using the best available amplitude bound.}
Near $t_*$ (where $t \asymp c^{1/2}$),
the sharpest pointwise bound is
$|\widehat{\Phi}_n^{(c)}(t)| \le C_\kappa c^{-1/2}$
(from Theorem~A or from (VI.2)
with $t \asymp c^{1/2}$).
This gives $f(\Delta,c)^2 = C_\kappa^2 c^{-1}$,
hence:
\[
\Delta \ge \frac{c_1}{C_\kappa^2 c^{-1}}
= c_1 C_\kappa^{-2}\,c.
\]
This is a lower bound of order $c$,
not $c^{-1/2}$ -- it says
$\Delta$ cannot be smaller than $O(c)$
based on mass alone,
which is a trivially weak statement.

\textbf{Using the crossover-zone amplitude.}
For $t \in \Zcross$,
Theorem~A gives the sharper bound
$|\widehat{\Phi}_n^{(c)}(t)|
\le C_\kappa c^{-1/2}|t-t_*|^{-1/4}$.
At the boundary $|t-t_*| = \Delta$:
$f = C_\kappa c^{-1/2}\Delta^{-1/4}$,
so $f^2 = C_\kappa^2 c^{-1}\Delta^{-1/2}$,
and the mass bound gives:
\[
2\Delta \cdot C_\kappa^2 c^{-1}\Delta^{-1/2}
= 2C_\kappa^2 c^{-1}\Delta^{1/2} < \tfrac{1}{2},
\]
i.e.\ $\Delta < c^2/(16 C_\kappa^4)$,
which is an \emph{upper} bound on $\Delta$,
not a lower bound.

\textbf{Conclusion.}
No combination of the inputs (a) and (b)
can produce a lower bound
$\Delta_\varepsilon \ge c_1 c^{-1/2}$,
because:
\begin{enumerate}[\rm(i)]
\item Using $|f| \le C_\kappa c^{-1/2}$
gives a lower bound $\Delta \ge O(c)$,
trivially weaker than $c^{-1/2}$.
\item Using $|f| \le C_\kappa c^{-1/2}|t-t_*|^{-1/4}$
gives an upper bound on $\Delta$,
not a lower bound.
\item The bottleneck is that
the dominant $L^2$-mass
($\tfrac{1}{2}$, bulk)
resides in the algebraic zone
$|t - t_*| \ge c^{-1/3}$,
far from $t_*$.
Any method that does not distinguish
pointwise non-cancellation
from cancellation in this zone
will fail to see the $c^{-1/2}$ scale.
\end{enumerate}
\end{proof}

\begin{remark}[Structural obstruction]
\label{rem:structural}
Proposition~\ref{prop:lower-barrier}
identifies the precise obstruction:
the lower bound requires knowing that
$|\widehat{\Phi}_n^{(c)}(t)|$ does
\emph{not} cancel on a set of measure
$\gtrsim c^{-1/2}$ near $t_*$.
Global mass arguments, by their nature,
cannot detect local non-cancellation:
they see only the total $L^2$-budget,
which is dominated by the algebraic zone
(Corollary~\ref{cor:zone-table-complete}).
The required non-cancellation statement
is Problem~\ref{prob:nonvanish},
which is the central open problem
for Paper~VIII.
\end{remark}

\begin{remark}[What would be needed]
\label{rem:lower-needed}
The correct lower bound requires
$|\widehat{\Phi}_n^{(c)}(t)| \ge c_\kappa c^{-1/4}t^{-1/4}$
on a set of measure $\gtrsim c^{-1/2}$
in the algebraic zone near $\tstar$;
see Problem~\ref{prob:nonvanish}.
\end{remark}

%%═══════════════════════════════════════════════════════════════════
\section{Open Problems}
\label{sec:problems}
%%═══════════════════════════════════════════════════════════════════

\begin{problem}[Non-cancellation lower bound]
\label{prob:nonvanish}
Prove that there exist $c_1(\kappa)>0$
and $E_c \subset \{|t-\tstar|\le c^{-1/2}\}$
with $|E_c| \ge c_1 c^{-1/2}$ such that
$|\widehat{\Phi}_n^{(c)}(t)| \ge c_1 c^{-1/2}$
for all $t\in E_c$, $n\le\kappa c$, $c\ge c_0$.
This is the central missing ingredient for
$\Delta_\varepsilon \ge c_1 c^{-1/2}$
(Paper~VIII).
\end{problem}

\begin{problem}[Sharp constant in Theorem~B]
\label{prob:sharp-const}
Determine
$C^*(\kappa,\varepsilon)
:= \limsup_{c\to\infty}
c^{1/2}\Delta_\varepsilon(c,\lfloor\kappa c\rfloor)$.
\end{problem}

%%═══════════════════════════════════════════════════════════════════
\begin{thebibliography}{99}

\bibitem{PaperV}
[Author],
\textit{Effective bandwidth and $L^2$-concentration
of prolate spheroidal wave functions: numerical
evidence for the $c^{-1/2}$ peak-width scaling},
Paper~V of this series, preprint, 2025.

\bibitem{PaperVI}
[Author],
\textit{Pointwise Mellin decay, log-free $L^2$ tails,
and the $c^{-2/3}$ upper bound for PSWFs},
Paper~VI of this series, preprint, 2025.

\bibitem{Slepian1961}
D.~Slepian and H.~O.~Pollak,
\textit{Prolate spheroidal wave functions,
Fourier analysis and uncertainty~I},
Bell System Tech.\ J.\ \textbf{40} (1961), 43--63.

\bibitem{Prolate1}
D.~Slepian,
\textit{Prolate spheroidal wave functions,
Fourier analysis and uncertainty~IV},
Bell System Tech.\ J.\ \textbf{43} (1964), 3009--3057.

\bibitem{Osipov2013}
A.~Osipov, V.~Rokhlin, and H.~Xiao,
\textit{Prolate Spheroidal Wave Functions
of Order Zero},
Springer, New York, 2013.

\bibitem{Olver1974}
F.~W.~J.~Olver,
\textit{Asymptotics and Special Functions},
Academic Press, New York, 1974.

\end{thebibliography}

\end{document}
